% SECCIÓN 2: LA RESPONSABILIDAD CIVIL MÉDICA EN EL ORDENAMIENTO JURÍDICO PERUANO
%
% Objetivos de esta sección:
%   1. Enmarcar la responsabilidad civil médica en el sistema general del CC de 1984:
%      - Régimen contractual (art. 1314 CC): obligación de medios vs. de resultado.
%      - Régimen extracontractual (art. 1969 CC): culpa y antijuridicidad.
%   2. Explicar la dualidad del régimen peruano y sus consecuencias prácticas
%      (plazo de prescripción, carga de la prueba, quantum indemnizatorio).
%   3. Identificar los cuatro elementos de la responsabilidad civil:
%      antijuridicidad, factor de atribución (culpa / riesgo), nexo causal, daño.
%   4. Situar la Ley 26842 (Ley General de Salud) y la Ley 29414
%      como normativa especial que complementa el CC.
%
% Referencias clave: Espinoza Espinoza, Taboada Córdova, De Trazegnies, art. 1969-1985 CC
% Extensión estimada: 1.200–1.800 palabras

\section{La responsabilidad civil médica en el ordenamiento jurídico peruano}
\label{sec:rc_medica}

